
\begin{table}[H]
    \centering
     \caption{Frecuencia de los factores clave o códigos detectados en las cinco metodologías activas.}
     
    \begin{tabular}{p{4.5cm}ccccc}
    \hline
       Factor	&	Aula Invertida	&	ABP	&	ABpr	&	AC	&	Gamificación	\\
     \hline
Capacitación docente	&	22	&	21	&	18	&	17	&	23	\\
Planificación del tiempo	&	23	&	18	&	20	&	21	&	22	\\
Comunicación previa de la planificación	&	23	&	0	&	0	&	0	&	0	\\
Retroalimentación al estudiante	&	23	&	21	&	18	&	23	&	24	\\
Selección de preguntas y problemas motivadoras y retadoras	&	0	&	24	&	21	&	22	&	0	\\
Diseño de actividades cooperativas	&	1	&	2	&	2	&	23	&	2	\\
Capacidad del estudiante para investigar	&	22	&	23	&	21	&	22	&	3	\\
Selección de herramientas tecnológicas	&	3	&	4	&	2	&	19	&	22	\\
 \hline
    \end{tabular}
    \caption*{\it{Nota}: Los números en cada celda indican el número de veces que el factor clave de implementación fue nombrado en los artículos analizados sobre cada una de las metodologías activas estudiadas. ABP: Aprendizaje basado en proyectos. ABpr: Aprendizaje basado en problemas. AC: Aprendizaje cooperativo. }
      \label{tab:FrecFactoresClave}
\end{table}



% Please add the following required packages to your document preamble:
% \usepackage{booktabs}
\begin{table}
\centering
\caption{\label{tab:BondadAjuste}Indicadores de bondad de ajuste del modelo de la dimensión 1, conocimiento y uso de herramientas tecnológicas para la innovación docente.}
\begin{tabular}{@{}lllllll@{}}
\toprule
\textbf{$\chi^{2}$} & \textbf{gl} & \textbf{p-valor} & \textbf{$\chi^{2}$/\textit{gl}} & \textit{TLI} & \textit{CFI} & \textit{RMSEA [IC]} \\ \midrule
7304255     & 860         & .000              & 8493           & .931         & .934         & .156 {[}.152-.159{]}   \\ \bottomrule
\end{tabular}
    
\end{table}
